% ==============================================================================
% 09_conclusion.tex
% Phụ trách: Thành viên 2 (Thực nghiệm, kết quả và kết luận)
% ==============================================================================

\section{Kết luận}
\label{sec:conclusion}

Nghiên cứu này đã thực hiện một cuộc đánh giá thực nghiệm toàn diện, nghiêm ngặt và có khả năng tái lập độc lập 100\% nhằm làm sáng tỏ tác động thực sự của ba chiến lược chunking---Fixed-size, Structure-aware và Prompt-based---lên toàn bộ chu trình RAG trên tài liệu kỹ thuật Markdown. Dưới thiết kế kiểm soát chặt chẽ với cùng một kho ngữ liệu Angular 384 tài liệu, cùng một mô hình embedding và generator, cùng giới hạn nghiêm ngặt của retrieval token budget 2.048 tokens và context budget 4.096 tokens, công trình đã giải đáp dứt khoát ba câu hỏi nghiên cứu cốt lõi:

\begin{itemize}
    \item \textbf{RQ1 (Retrieval Ranking và Early Precision):} Chiến lược Fixed-size chiếm ưu thế tuyệt đối ở các thứ hạng sớm với $\text{Hit@1} = 0.8071$ và $\text{MRR} = 0.8732$ nhờ kích thước chunk lớn ($447.1$ tokens) và cơ chế sliding-window overlap 64 tokens tạo ra mật độ ngữ nghĩa cục bộ dày đặc. Tuy nhiên, khi mở rộng phạm vi ứng viên, Structure-aware nhanh chóng bứt phá và xác lập kỷ lục dẫn đầu với $\text{Hit@10} = 1.0000$ và $\text{Recall@10} = 0.9857$ nhờ việc tôn trọng ranh giới heading hierarchy của AST parser và ngăn ngừa hiện tượng semantic dilution.
    
    \item \textbf{RQ2 (Khả năng Thu hồi Bằng chứng trong Ngân sách Token):} Dưới giới hạn retrieval token budget 2.048 tokens, Structure-aware thể hiện sức mạnh áp đảo toàn diện với $\text{Evidence Coverage} = 0.7476$ ($74.76\%$) và $\text{All-Evidence-Retrieved Rate} = 63.57\%$, vượt xa Fixed-size ($69.05\%$ và $57.14\%$). Khả năng đóng gói token tối ưu (budget packing efficiency đạt $99.57\%$) cho phép thuật toán greedy accumulator nhồi trung bình $14.64$ chunks hạt mịn độc lập vào context window mà không lãng phí tài nguyên vào phần văn bản trùng lặp thừa thãi.
    
    \item \textbf{RQ3 (Chất lượng Câu trả lời Sinh ra):} Fixed-size dẫn đầu về tổng thể Answer Token F1 ($0.3830$) và Token Precision ($0.2858$) do cung cấp ít khối văn bản hơn ($4.94$ chunks), bảo tồn dòng chảy diễn ngôn liền mạch và giảm thiểu gánh nặng tổng hợp (synthesis burden) cho downstream generator (\texttt{gpt-5-mini}); trong khi Prompt-based đạt Token Recall cao nhất ($0.7555$). Dù vậy, các kiểm định hoán vị có dấu (paired sign-flip permutation tests, Holm-adjusted $p > 0.82$) xác nhận rằng sự chênh lệch Token F1 giữa ba phương pháp chưa đạt mức có ý nghĩa thống kê, chứng minh rằng ưu thế thu hồi bằng chứng ở tầng retrieval không tự động chuyển hóa thành bước nhảy vọt ở tầng downstream generation.
\end{itemize}

Thông điệp cốt lõi được đúc kết từ công trình này là: \textit{Không tồn tại một chiến lược chunking tối ưu vạn năng cho mọi hệ thống RAG (no one-size-fits-all)}. Quyết định lựa chọn kiến trúc phải căn cứ chặt chẽ vào mục tiêu nghiệp vụ cụ thể: các kỹ sư RAG nên lựa chọn \textbf{Fixed-size} khi xây dựng các ứng dụng factoid QA hoặc tra cứu định nghĩa đơn giản ưu tiên độ trễ cực thấp và độ chuẩn xác ở Top-1; và bắt buộc phải triển khai \textbf{Structure-aware} khi đối mặt với các bài toán kiểm toán tuân thủ, tài liệu API chuyên sâu hoặc tổng hợp đa bước (multi-hop reasoning) nơi việc thu hồi đầy đủ 100\% bằng chứng kỹ thuật là điều kiện tiên quyết. Các phát hiện về hiện tượng Context Fragmentation trong nghiên cứu này mở ra triển vọng ứng dụng các kiến trúc lai (hybrid approaches) như Parent-Document Retrieval và Dynamic Chunk Merging nhằm dung hòa trọn vẹn giữa độ bao phủ bằng chứng của bộ parser cấu trúc và tính liền mạch ngữ cảnh cho các mô hình ngôn ngữ lớn thế hệ mới.
